\documentclass[10pt]{amsart}
\usepackage[a4paper,margin=23mm]{geometry}
\usepackage{amsmath,amssymb,amsthm,mathtools,hyperref,graphicx,xspace,letltxmacro,xpatch}
\usepackage[english]{babel}
\usepackage{xfrac,xcolor}

\usepackage[scr=rsfs,cal=euler]{mathalfa}
\multlinegap0pt

\newcommand{\Changel}{\mathcode`l="7160}
\newcommand{\Changelback}{\mathcode`l="716C}
\newcommand{\NoChangel}[1]{%
\expandafter\let\csname old\string#1\endcsname=#1
\let#1=\relax
\newcommand{#1}{\mathop{\mathcode`l="716C\csname old\string#1\endcsname\mathcode`l="7160}\displaylimits}%
}

\NoChangel{\log}
\NoChangel{\ln}
\NoChangel{\lim}
\NoChangel{\limsup}
\NoChangel{\liminf}
\NoChangel{\varlimsup}
\NoChangel{\varliminf}
\NoChangel{\varinjlim}
\NoChangel{\varprojlim}

\newcommand\mto{\mathchoice{\longmapsto}{\mapsto}{\mapsto}{\mapsto}}
\newcommand\To[1]{\mathchoice{\xrightarrow{\textstyle\kern4pt#1\kern3pt}}{\stackrel{#1}{\longrightarrow}}{}{}}
\newcommand{\Psfrac}[2]{(\sfrac{#1}{#2})}
\newcommand{\psfrac}[2]{\sfrac{(#1)}{#2}}
\newcommand{\spfrac}[2]{\sfrac{#1}{(#2)}}
\newcommand{\pspfrac}[2]{\sfrac{(#1)}{(#2)}}

\newcommand{\simm}{\mathop{\sim}\displaylimits}

\usepackage[nameinlink,capitalise]{cleveref}
\crefdefaultlabelformat{#2\textup{#1}#3}
\crefname{figure}{Figure}{Figures}
\crefname{equation}{Equation}{Equations}


\theoremstyle{plain}
\newtheorem{thm}{Theorem}[subsection]
\newtheorem{lem}[thm]{Lemma}
\newtheorem{coro}[thm]{Corollary}
\newtheorem{thmalph}{Theorem}
\newtheorem{coralph}[thmalph]{Corollary}
\renewcommand{\thethmalph}{\Alph{thmalph}}
\newtheorem{prop}[thm]{Proposition}
\newtheorem{as}[thm]{Assumption}
\newtheorem{con}[thm]{Condition}
\newtheorem{As}[thm]{Assumption}
\newtheorem{conj}{Conjectural equality}
\newtheorem{theorem}[thm]{Theorem}
\newtheorem{corollary}{Corollary}
\newtheorem*{claim}{Claim}

\theoremstyle{definition}
\newtheorem{rmrk}[thm]{Remark}
\newtheorem{example}[thm]{Example}
\theoremstyle{remark}
\newtheorem*{ack}{Acknowledgements}
\newtheorem{defi}[thm]{Definition}
\newtheorem{ex}[thm]{Example}
\newtheorem{setting}[thm]{Setting}

\newcommand{\R}{\mathbb{R}}
\newcommand{\Z}{\mathbb{Z}}
\newcommand{\N}{\mathbb{N}}
\newcommand{\K}{\mathbb{K}}
\renewcommand{\H}{\mathbb{H}}
\newcommand{\C}{\mathbb{C}}
\newcommand{\matrx}{\matrix}
\newcommand{\abs}[1]{\left|#1\right|}
\newcommand{\eps}{\varepsilon}
\newcommand {\red}{\color[rgb]{1,0,0}}
\newcommand {\blue}{\color[rgb]{0,0,1}}
\newcommand{\A}{\mathbb{A}}
\newcommand{\Q}{\mathbb{Q}}
\newcommand{\ii}{\mathrm i}
\newcommand*{\rom}[1]{\romannumeral}
\renewcommand{\Re}{\textup{Re}}
\renewcommand{\Im}{\textup{Im}}
\newcommand{\rot}{\mathrm{rot}}
\newcommand{\geod}{\mathrm{geod}}
\newcommand{\bma}{\begin{pmatrix}}
\newcommand{\ema}{\end{pmatrix}}
\newcommand{\bigzero}{\mbox{\normalfont\large\bfseries 0}}
\newcommand{\bigId}{\mbox{\normalfont\large\bfseries \textrm{I}}}
\DeclareMathOperator{\fl}{flat}
\DeclareMathOperator{\SL}{SL}
\let\sl\relax
\DeclareMathOperator{\sl}{\mathfrak{sl}}
\DeclareMathOperator{\PSL}{PSL}
\DeclareMathOperator{\GL}{GL}
\DeclareMathOperator{\SO}{SO}
\DeclareMathOperator{\PSO}{PSO}
\DeclareMathOperator{\SU}{SU}
\DeclareMathOperator{\U}{U}
\DeclareMathOperator{\Spin}{Spin}
\DeclareMathOperator{\IM}{Im}
\DeclareMathOperator{\RE}{Re}
\DeclareMathOperator{\Tr}{Tr}
\DeclareMathOperator{\tr}{tr}
\DeclareMathOperator{\End}{End}
\DeclareMathOperator{\Vol}{Vol}
\DeclareMathOperator{\rank}{rank}
\DeclareMathOperator{\Id}{Id}
\DeclareMathOperator{\Ad}{Ad}
\DeclareMathOperator{\ad}{ad}
\DeclareMathOperator{\T}{T}
\DeclareMathOperator{\Norm}{Norm}
\DeclareMathOperator{\cent}{Centr}
\DeclareMathOperator{\spec}{spec}
\DeclareMathOperator{\diam}{diam}
\DeclareMathOperator{\order}{order}
\DeclareMathOperator{\sign}{sign}
\DeclareMathOperator{\Res}{Res}
\DeclareMathOperator{\prim}{prime}
\DeclareMathOperator{\supp}{supp}
\DeclareMathOperator{\rk}{rk}
\DeclareMathOperator{\area}{area}
\DeclareMathOperator{\ind}{ind}
\DeclareMathOperator{\Ind}{Ind}
\DeclareMathOperator{\coker}{coker}
\DeclareMathOperator{\Cl}{Cl}
\DeclareMathOperator{\AS}{AS}
\DeclareMathOperator{\dom}{dom}
\DeclareMathOperator{\gr}{gr}
\DeclareMathOperator{\Rep}{Rep}
\DeclareMathOperator{\ev}{even}
\DeclareMathOperator{\im}{Im}
\DeclareMathOperator{\Ker}{Ker}
\DeclareMathOperator{\Flat}{Flat}
\DeclareMathOperator{\Mat}{Mat}
\DeclareMathOperator{\RHS}{RHS}
\DeclareMathOperator{\LHS}{LHS}
\DeclareMathOperator{\Cent}{Cent}
\DeclareMathOperator{\Fix}{Fix}
\numberwithin{equation}{section}
\DeclareMathOperator{\dive}{div}
\DeclareMathOperator{\grad}{grad}
\DeclareMathOperator{\Arg}{Arg}
\DeclareMathOperator{\tor}{tor}
\DeclareMathOperator{\mult}{mult}

\renewcommand{\thethm}{%
\ifnum\value{subsection}=0
\thesection
\else
\thesubsection
\fi
.\arabic{thm}%
}

\let\dpl\displaystyle
\providecommand{\skpt}{\par\smallskip\noindent}
\providecommand{\pointir}{.\ }
\providecommand{\eg}{e.g.}\providecommand{\ie}{i.e.}\providecommand{\etc}{etc.}\providecommand{\resp}{resp.}\providecommand{\cf}{cf.}
\providecommand{\mainmatter}{}\providecommand{\backmatter}{}
\let\Subsection\subsection\let\Subsubsection\subsubsection
\AtBeginDocument{\gdef\shorttitle{Twisted Ruelle zeta and Reidemeister-Turaev torsion}}
\setlength{\emergencystretch}{2em}
